% !TEX program = lualatex
% Generated by new_lecture_note.sh
\documentclass[11pt]{article}

\usepackage{fontspec}
\usepackage{microtype}
\usepackage{mathtools,amssymb,amsfonts,amsthm,bm}
\usepackage{enumitem}
\usepackage{booktabs,array,xcolor}
\usepackage{geometry,fancyhdr,setspace}
\usepackage[hidelinks]{hyperref}
\usepackage[nameinlink,noabbrev]{cleveref}

\geometry{margin=1in,headheight=15pt}
\setstretch{1.12}
\setlength{\parindent}{0pt}
\setlength{\parskip}{0.55em plus 0.12em minus 0.08em}
\setlength{\abovedisplayskip}{0.8em plus 0.2em minus 0.1em}
\setlength{\belowdisplayskip}{0.8em plus 0.2em minus 0.1em}
\allowdisplaybreaks[3]
\setlist{leftmargin=*,topsep=0.35em,itemsep=0.15em,parsep=0pt}

% ---------- Metadata ----------
\newcommand{\studentname}{Keshav Ramamurthy}
\newcommand{\coursename}{Math 113}
\newcommand{\courselongname}{Abstract Algebra}
\newcommand{\lecturenumber}{03}
\newcommand{\lecturetitle}{Lecture 03}
\newcommand{\lecturedate}{\today}

% ---------- Reusable notation ----------
\newcommand{\N}{\mathbb{N}}
\newcommand{\Z}{\mathbb{Z}}
\newcommand{\Q}{\mathbb{Q}}
\newcommand{\R}{\mathbb{R}}
\newcommand{\C}{\mathbb{C}}
\newcommand{\F}{\mathbb{F}}
\newcommand{\K}{\mathbb{K}}
\newcommand{\E}{\mathbb{E}}
\newcommand{\Prob}{\mathbb{P}}
\newcommand{\1}{\mathbf{1}}
\newcommand{\eps}{\varepsilon}
\newcommand{\dd}{,\mathrm{d}}
\newcommand{\abs}[1]{\left\lvert#1\right\rvert}
\newcommand{\norm}[1]{\left\lVert#1\right\rVert}
\newcommand{\ip}[2]{\left\langle#1,#2\right\rangle}
\newcommand{\set}[1]{\left{#1\right}}
\newcommand{\given}{,\middle\vert,}
\newcommand{\indep}{\perp!!!\perp}
\newcommand{\lto}{\longrightarrow}
\newcommand{\deriv}[2]{\frac{\mathrm{d}#1}{\mathrm{d}#2}}
\newcommand{\pderiv}[2]{\frac{\partial#1}{\partial#2}}
\DeclareMathOperator{\Aut}{Aut}
\DeclareMathOperator{\Hom}{Hom}
\DeclareMathOperator{\Ker}{ker}
\DeclareMathOperator{\im}{im}
\DeclareMathOperator{\rank}{rank}
\DeclareMathOperator{\nullity}{nullity}
\DeclareMathOperator{\Span}{span}
\DeclareMathOperator{\Var}{Var}
\DeclareMathOperator{\Cov}{Cov}

% ---------- Note environments ----------
\newtheoremstyle{notestyle}{0.7em}{0.7em}{\itshape}{}{}{\bfseries}{.5em}{}
\theoremstyle{notestyle}
\newtheorem{theorem}{Theorem}[section]
\newtheorem{lemma}[theorem]{Lemma}
\newtheorem{proposition}[theorem]{Proposition}
\newtheorem{corollary}[theorem]{Corollary}
\theoremstyle{definition}
\newtheorem{definition}[theorem]{Definition}
\newtheorem{example}[theorem]{Example}
\newtheorem{remark}[theorem]{Remark}
\newenvironment{claim}{\par\noindent\textbf{Claim.}\ }{\par}
\newenvironment{idea}{\par\noindent\textit{Idea.}\ }{\par}
\newenvironment{proofsketch}{\par\noindent\textbf{Proof sketch.}\ }{\par}

% ---------- Header ----------
\pagestyle{fancy}
\fancyhf{}
\fancyhead[L]{\coursename}
\fancyhead[C]{\lecturetitle}
\fancyhead[R]{\studentname}
\fancyfoot[C]{\thepage}

\title{\vspace{-1.5em}\textbf{\coursename\ -- \lecturetitle}}
\author{\studentname}
\date{\lecturedate}

\begin{document}
\maketitle

\section{Subgroups (Section 5)}
\subsection{Basic Relationships between Subgroups}

\begin{definition}[Subgroup]
\label{def:subgroup}
A subgroup $H$ of a group $G$ is a subset, closed under the operation $*$ in $G$, such
that $\langle H, *\rangle$ is a group. Write

$$
H \le G \qquad (H \subseteq G).
$$

\end{definition}

\begin{example}[Subgroup Examples]
\label{ex:subgroup}
Take $H = (\mathbb{Z}, +)$ and $G = (\mathbb{Q}, +)$.
\end{example}

\begin{example}[\textbf{NOT A GROUP}]
Taking:

$$
H = (\mathbb{Z}, \cdot) \qquad \text{and} \qquad G = (\mathbb{Q}, +)
$$

Must use induced operation $(+)$ from $G$.
\end{example}

\begin{example}[Smallest Subgroup of $G$]
\label{ex:trivial-subgroup}
Take ${e}$. \par
It's closed as $e * e = e$, in fact it is the trivial group.
\end{example}

\begin{example}[$G$ is a subgroup of itself]
\label{ex:self-subgroup}
$G \subseteq G.$
\end{example}

\begin{theorem}[\textbf{Groups inside other groups}]
\label{thm:subgroup-transitivity}
If $K \le H$ and $H \le G$, then $K \le G$.
\end{theorem}

\begin{proof}
Since $K \subseteq H \subseteq G$, we know that $K \subseteq G$.\par
Since $K \le H$, we know that $K$ is closed under $*$, and
$\langle K, *\rangle$ is a group, so $K \le G$.
\end{proof}

\subsection{Order}

\begin{definition}[Order of a group]
\label{def:group-order}
For a finite group $G$, we call $|G|$ the number of elements in $G$, its order.
\end{definition}

\newpage

\begin{remark}[Groups of Order $2$ and $3$]
\label{rem:orders-two-three}
These are the unique groups of order $2$ and $3$.
\end{remark}

$$
\begin{array}{c|cc}
 & e & a \\ \hline
e & e & a \\
a & a & e
\end{array}
\qquad
\begin{array}{c|ccc}
 & e & a & b \\ \hline
e & e & a & b \\
a & a & b & e \\
b & b & e & a
\end{array}
$$

$$
\boxed{G_2 = \{e,a\}}
\qquad\qquad
\boxed{G_3 = \{e,a,b\}}
$$

There are no trivial subgroups.

\begin{example}[Groups of Order 4]
\label{ex:order-four}

$$
\boxed{V_4 = \{e,a,b,c\}}
$$

$$
\begin{array}{c|cccc}
\cdot & e & a & b & c \\ \hline
e & e & a & b & c \\
a & a & e & c & b \\
b & b & c & e & a \\
c & c & b & a & e
\end{array}
$$

What are its subgroups?

\par
We will see: If $G$ is finite and $H \le G$, then $|H|$ divides $|G|$.

\begin{theorem}[Checking for Subgroups]
\label{thm:subgroup-test}
A subset $H$ of $G$ is a subgroup if
\begin{itemize}
\item $H$ is closed under the operation of $G$;
\item the identity element $e$ of $G$ is in $H$;
\item for any $a \in H$, the inverse $a^{-1}$ is also in $H$.
\end{itemize}
\end{theorem}

\begin{proof}
H needs an identity to be a subgroup, and that identity element must be $e$, the same as
$G$ because otherwise

$$
e_H * h = h
$$

would not imply that the identity of $H$ is the identity of $G$.

Associativity works here because the induced operation is associative.

The inverse of $a \in H$ is $a^{-1}$. Just prove this by assuming that $a$ has an alternate
inverse $a_H$.
\end{proof}

\end{example}

\begin{definition}[General Linear Group]
\label{def:general-linear-group}
The General Linear Group $GL_n(\mathbb{C})$ is the set of $n \times n$ invertible matrices
with entries in $\mathbb{C}$.
\end{definition}
\subsection{Linear Subgroups}
\begin{definition}[Simple Linear Group]
\label{def:special-linear-group}
Define the Special Linear Group $SL_n(\mathbb{C})$ as the set of $n \times n$ invertible
matrices with entries in $\mathbb{C}$ with determinant $1$.
\end{definition}
\newpage
\begin{proposition}[subgroup]
    \label{def:label}
    We claim the Simple Linear Group is a subgroup of $GL_n$
Proposition
\end{proposition}
\begin{proof}
    We see immediately that $SL_n{\mathbb{C}}\subseteq GL_n(\mathbb{C})$, since $\det
    \ne 0\implies$ invertible
    \par Then we see that if $\det(A) = \det(B) = 1 \implies \det(AB) = \det(A)
    \det(B) = 1$, establishing closure
    \par We see the identity $I$ where we see $\det I =1$
    \par Lastly, we look at the inverse, where we see $\det{A^{-1}}=\frac{1}{\det A} =
    1$, so it's closed under inverses
\end{proof}
\subsection{Cyclic Subgroups}
We have for $a \in G$, let $a^n = a * a * \cdots * a$ n times. Where $a^{-1} = a^{1}$,
$a^{-n} = (a^{-1})^{n}.$ Let $a^{0} = e.$
\begin{theorem}[Cyclic Subgroups]
\label{thm:label}
Let $a$ be an element  of a group $G$. Then $$H = \{a^n | n \in \mathbb{Z}\}$$ is a
subgroup  of $G$, and every subgroup of $G$ which contains $a$ must contain $H.$
\end{theorem}
\begin{proof}
    We see that $a^{n} * a^{m} = a^{n + m}$ pretty easily from associativity.
    \par Then $e = a^{0} \in H.$
    \par Lastly $(a^n)^{-1} = a^{-n} \in H$
\end{proof}
\begin{definition}[Group "Generation"]
\label{def:label}
An element $a \in G$ generates $G$ is $\langle a\rangle  = G.$\par  G is cyclic if there is
some element $a$ that generates it.
\end{definition}
\section{Definitions and notation}

\begin{definition}[First definition]
Write the definition and one sentence explaining why it matters.
\end{definition}

\section{Claims, theorems, and proof sketches}

\begin{claim}
State the claim in one precise sentence.
\end{claim}

\begin{proofsketch}
Record the key idea, the first useful equation, and the step that still needs checking.
\end{proofsketch}

\section{Examples and calculations}

\begin{example}[Lecture example]
Write the setup, the calculation, and the conclusion.
\end{example}

\section{Questions and follow-up}

\begin{itemize}
\item What is still unclear?
\item Which exercise or reading should I revisit?
\item TODO:
\end{itemize}

\end{document}

